%%%%%%%%%%%%%%%%%%%%%%%%%%%%%%%%%%%%%%%%%%%%%%%%%%%%%%%%%%%%%%%%%%%%%%%%%%%%%%%%
%2345678901234567890123456789012345678901234567890123456789012345678901234567890
%        1         2         3         4         5         6         7         8

\documentclass[letterpaper, 10 pt, conference]{ieeeconf}  % Comment this line out if you need a4paper

%\documentclass[a4paper, 10pt, conference]{ieeeconf}      % Use this line for a4 paper

\IEEEoverridecommandlockouts                              % This command is only needed if 
                                                          % you want to use the \thanks command

\overrideIEEEmargins                                      % Needed to meet printer requirements.
        % 匯入你的 .bib 檔名 (不需加 .bib 副檔名)
%In case you encounter the following error:
%Error 1010 The PDF file may be corrupt (unable to open PDF file) OR
%Error 1000 An error occurred while parsing a contents stream. Unable to analyze the PDF file.
%This is a known problem with pdfLaTeX conversion filter. The file cannot be opened with acrobat reader
%Please use one of the alternatives below to circumvent this error by uncommenting one or the other
%\pdfobjcompresslevel=0
%\pdfminorversion=4

% See the \addtolength command later in the file to balance the column lengths
% on the last page of the document

% The following packages can be found on http:\\www.ctan.org
%\usepackage{graphics} % for pdf, bitmapped graphics files
%\usepackage{epsfig} % for postscript graphics files
%\usepackage{mathptmx} % assumes new font selection scheme installed
%\usepackage{times} % assumes new font selection scheme installed
%\usepackage{amsmath} % assumes amsmath package installed
%\usepackage{amssymb}  % assumes amsmath package installed
\usepackage{cite}
\title{\LARGE \bf
Physics-Informed Deep Stereo: Active Cross-Polarization 
for Data-Efficient Transparent Obstacle Avoidance
}


\author{Lin. Po-Ting 
\thanks{*This work was not supported by any organization}% <-this % stops a space
\thanks{$^{1}$Albert Author is with Faculty of Electrical Engineering, Mathematics and Computer Science,
        University of Twente, 7500 AE Enschede, The Netherlands
        {\tt\small albert.author@papercept.net}}%
\thanks{$^{2}$Bernard D. Researcheris with the Department of Electrical Engineering, Wright State University,
        Dayton, OH 45435, USA
        {\tt\small b.d.researcher@ieee.org}}%
}


\begin{document}



\maketitle
\thispagestyle{empty}
\pagestyle{empty}


%%%%%%%%%%%%%%%%%%%%%%%%%%%%%%%%%%%%%%%%%%%%%%%%%%%%%%%%%%%%%%%%%%%%%%%%%%%%%%%%
\begin{abstract}

This electronic document is a ÒliveÓ template. The various components of your paper [title, text, heads, etc.] are already defined on the style sheet, as illustrated by the portions given in this document.

\end{abstract}


%%%%%%%%%%%%%%%%%%%%%%%%%%%%%%%%%%%%%%%%%%%%%%%%%%%%%%%%%%%%%%%%%%%%%%%%%%%%%%%%
\section{INTRODUCTION}

Transparent obstacles are difficult to detect using traditional methods. Their characteristics, including light transmission, spectral reflection, and refraction distortion, pose significant challenges for detection.

Traditional obstacle detection methods struggle with transparent obstacles. 
LiDAR sensors fail due to light transmission through transparent surfaces 
\cite{tibebu_lidar-based_2021}. RGB-D cameras produce incomplete depth maps with large missing regions due to spectral reflection 
and refraction distortion \cite{lysenkov_recognition_2013}. Although optical transparency does not affect 
ultrasonic sensors, they lack the spatial resolution and range required for dense 3D reconstruction in robotic navigation \cite{borenstein_vector_1991}.

Recent learning-based methods have shown promise in addressing this challenge. 
Depth completion approaches train neural networks to infer complete depth maps 
from incomplete RGB-D inputs \cite{sajjan_clear_2020}.
However, these methods are data-intensive and require tens of thousands of 
synthetic training images to achieve satisfactory performance. Active stereo methods \cite{shi_asgrasp_2024} have been proposed to improve depth estimation 
for transparent objects by leveraging structured light patterns. While these 
approaches reduce the dependency on scene texture, they still require extensive 
training data to handle the unique optical properties of transparent materials, 
such as refraction and specular reflection.

To address this limitation, we propose a physics-informed deep stereo 
framework for data-efficient transparent obstacle detection. Our approach 
leverages the physical principle that transparent objects produce highly 
polarized specular reflections. By capturing stereo pairs at orthogonal 
polarization angles (0° and 90°), the intensity difference reveals transparent 
object boundaries with high contrast, independent of texture or ambient lighting. 
This physics-based input representation enables fine-tuning of existing stereo 
networks with minimal training data, eliminating the need for large-scale 
specialized datasets.

\section{RELATED WORK}

\subsection{Hardware-Enhanced Approaches for Transparent Object Detection}
Modern transparent object detection relies primarily on deep learning-based 
visual models. A portion of recent research has focused on optimizing input 
representations to improve data efficiency, enabling networks to learn from 
fewer training samples by providing more informative inputs.

Thermal imaging methods exploit the fact that transparent objects have 
different thermal properties than their surroundings, making them visible 
in infrared images regardless of optical transparency 
\cite{huo_glass_2023}. However, thermal imaging 
requires temperature differences between the object and background, limiting 
its applicability in environments with uniform temperatures. Additionally, 
thermal cameras are significantly more expensive than RGB cameras, hindering 
widespread deployment in robotic systems.


\subsection{Polarization-Based Vision}
\subsection{Our Approach}

\section{METHOD}
\section{EXPERIMENTS}
\section{RESULTS}
\section{CONCLUSIONS}
\section*{APPENDIX}
\bibliographystyle{IEEEtran.bst}  % 指定使用 IEEE 標準格式
\bibliography{IEEEfull.bib}










\end{document}
